\scanpage{18}
triviale, \deleted{on en conclut} d'où
\[
A'=A+\mathfrak m'
\]
on aura
\[
\mathfrak m A'=\mathfrak m A+\mathfrak m\mathfrak m'\subset\mathfrak m+\mathfrak m'{}^2
\]
et compte tenu de $(*)$ on aura
\[
\mathfrak m A'=\mathfrak m,
\]
C'est à dire on est encore sous les
conditions du N° 1, prop. 3. On conclut
comme dessus (avec les notations du N° 1)
% slanted
\margin{\uline{\unsure{Référence|vérifier} au Lemme préliminaire}}
que la suite
\diagram{\begin{tikzcd}[ampersand replacement=\&,row sep=tiny]
A_0 \arrow[r] \& A'_0 \arrow[r,shift left] \arrow[r,shift right] \& A''_0\\
k \arrow[u,equal] \& \&
\end{tikzcd}}
est encore exacte, et satisfait à \deleted{\unsure{une|encore}} la
\unsure{$\overline{\mathrm m}$|même} condition que la suite du départ :
$x'_0\in A'_0-A_0$ implique que $\delta(x'_0)$ engendre \deleted{$\unsure{\Delta_{A_0}|\Delta}$} $\Delta_0$.

On peut donc encore supposer que $A$ est un corps
$k$, et de plus $A'$ est local à idéal maximal
de carré nul, donc de la forme
\[
A'\simeq D_k(V)\simeq k+V
\]
où $V$ est un vectoriel de dimension finie sur $k$.
% archive sheet number: 17
